%------------------------------------------------------------------------------%
\begin{question}
  Encontre e esboce as curvas que representam o corte da superfície
  \[
    4x^2 + y^2 = 16
  \]
  pelos planos \(z=0\) e \(y=0\)
\end{question}

%------------------------------------------------------------------------------%
\begin{resolution}{Solução}
  No plano \(z=0\) a equação
  \[
    4x^2+y^2=16
  \]
  \pause
  permanece inalterada,
  \pause
  porém, agora representa a elipse
  \[
    \frac{x^2}{2^2} + \frac{y^2}{4^2} = 1
  \]
  \pause
  que intercepta o eixo $x$ nos pontos \((2, 0)\) e \((-2, 0)\)

  e o eixo $y$ nos pontos \((4, 0)\) e \((-4, 0)\)
\end{resolution}

%------------------------------------------------------------------------------%
\begin{resolution}{Solução}
  \onslide<+->{No plano \(y=0\) a equação
  \[
    4x^2+y^2=16
  \]}
  \onslide<+->{se reduz a}
  \begin{align*}
    \onslide<+->{4x^2 + 0 & = 16    \\}
    \onslide<+->{x^2      & = 4     \\}
    \onslide<+->{\abs{x}  & = 2     \\}
    \onslide<+->{x        & = \pm 2}
  \end{align*}
  \onslide<+->{que representa as retas verticais \(x=2\) e \(x=-2\)}
\end{resolution}

%------------------------------------------------------------------------------%
\begin{resolution}{Solução}
\begin{columns}
\column{0.5\linewidth}
  \centering
  \(z=0\) \\
  \medskip
  \includegraphics{C_04_exemplo_1_z0}
\column{0.5\linewidth}
  \centering
  \(y=0\) \\
  \medskip
  \includegraphics{C_04_exemplo_1_y0}
\end{columns}
\end{resolution}

%------------------------------------------------------------------------------%
